\subsection*{The 'with-in' Block and 'where' Clause}

Similar to 'parallel-by', 'with-in' statement can also bind \emph{i-tuples} to a \emph{mdspan}. The below code shows an example.

\begin{lstlisting}[language=croqtile,caption={The with-in Block and where Clause}]
with index in [10, 10] {
  // index is ituple with 2 elements
}
\end{lstlisting}

Here, 'index' is a \emph{i-tuple} with 2 elements. Sometimes programmers prefer that the 2 elements being named. This is possible by using the below syntax.

\begin{lstlisting}[language=croqtile,caption={The with-in Block and where Clause}]
with {x, y} in [10, 10] {...}
\end{lstlisting}

Or even to name both the \emph{i-tuple} and its elements.

\begin{lstlisting}[language=croqtile,caption={The with-in Block and where Clause}]
with index = {x, y} in [10, 10] {...}
\end{lstlisting}

You name 'index' as a \textbf{bounded ituple} in such scenarios.

While 'with-in' may seem like 'parallel-by', it operates differently. One significant difference is 'with-in' statement does not have implication for parallelism. The code block inside 'with-in' statement is sequentially executed. It does nothing more than creating the \emph{bounded-ituple}.

In addition, programmers could append a 'where' clause. For example,

\begin{lstlisting}[language=croqtile,caption={The with-in Block and where Clause}]
with {m, n} in [M, N], {n_p, k} in [N_P, K] where n_p <-> n {
  // matmul implements with m,n,K. n_p is no long useful.
}
\end{lstlisting}

The code snippet requires 'n' and 'n\_p' to have an identical value in all iterations. Thus inside the 'with-in' block, it is possible to replace'n' whenever 'n\_p' is required, or the opposite. Such a facility is useful to program many AI kernels. Programmers should use operation '\textless{}-\textgreater{}' to establish such relations.
